% Section V — Results
% Draft 2026-10-20. Sources: eval/results/retrieval_*.json, benchmark_*.json,
% faithfulness_*.json; summarised in eval/ANSWER_EVAL_REPORT.md.

\section{Results}
\label{sec:results}

\subsection{Retrieval: Graph versus Dense}

\begin{table}[t]
\caption{Retrieval quality, nDCG@8. Graph retrieval in oracle mode; best of graph and \texttt{mxbai} in bold.}
\label{tab:retrieval}
\centering
\setlength{\tabcolsep}{4pt}
\begin{tabular}{lrccc}
\hline
Question type & $n$ & Graph & Dense (mxbai) & Dense (nomic) \\
\hline
Dataset search       & 8 & \textbf{0.817} & 0.688 & 0.259 \\
Method comparison    & 8 & \textbf{0.660} & 0.535 & 0.166 \\
Claim verification   & 8 & \textbf{0.657} & 0.563 & 0.340 \\
Method usage         & 9 & 0.710 & \textbf{0.817} & 0.532 \\
Paper lookup         & 8 & 0.583 & \textbf{0.724} & 0.247 \\
Topic search         & 9 & 0.850 & \textbf{0.979} & 0.866 \\
\hline
All                  & 50 & 0.716 & 0.725 & 0.413 \\
Test split           & 25 & 0.754 & 0.712 & 0.435 \\
\hline
\multicolumn{5}{p{0.93\columnwidth}}{\footnotesize Test split: graph end-to-end (\texttt{qwen2.5:7b}) 0.695; RRF of graph and mxbai 0.792.}
\end{tabular}
\end{table}

Table~\ref{tab:retrieval} compares the two retrievers. Overall they are on par:
0.716 versus 0.725 nDCG@8 on all questions, with a paired bootstrap interval of
the difference of $[-0.10, 0.09]$. During development, graph retrieval improved
from 0.527 to 0.754 on the test split, measured on the same final gold set, so
the parity is with a tuned graph retriever.

The per-type results are not on par. Graph retrieval is better where the
question has structure that maps onto graph relations: datasets used by a
paper, papers reporting both compared methods, and benchmark results that
decide a claim. Dense retrieval is better on broad topics, on method usage and on finding a
paper by its title, where textual similarity is sufficient and graph retrieval
depends on how an entity is spelled in ORKG.

The two systems retrieve largely different papers: on average only 0.42 of the
8 papers are shared. A reciprocal rank fusion of the two lists reaches 0.792 on
the test split, above both. In end-to-end mode graph retrieval loses 0.06
(0.754 to 0.695) through misclassified questions and wrongly extracted entities,
a source of error that dense retrieval does not have. Finally, the dense result
depends strongly on the embedder: with \texttt{nomic-embed-text} nDCG@8 drops to
0.413, and to 0.166 on method comparisons. The \texttt{mxbai} baseline is thus a
strong one.

\subsection{Answers: Graph versus Dense}

\begin{table}[t]
\caption{Answers with the same LLM and only the retriever changed (all 50 questions). E2E: strict end-to-end score; Cl.: grade-2 papers cited per claim question; MC: E2E on method comparison. For \texttt{qwen2.5:7b}, run 1 / run 2. Faithful for \texttt{gpt-oss-20b} on the 49 answers judged with both retrievers.}
\label{tab:answers}
\centering
\setlength{\tabcolsep}{2.5pt}
\begin{tabular}{llccccc}
\hline
Model & Retr. & E2E & Score & Faithful & Cl. & MC \\
\hline
qwen2.5 & graph & .748/.696 & .758/.726 & .771/.758 & 2.25/2.38 & .64/.60 \\
        & dense & .731/.733 & .809/.811 & .739/.733 & 0.38/1.25 & .37/.51 \\
gpt-oss-20b & graph & .745 & .830 & .688 & 2.12 & .68 \\
            & dense & .745 & .841 & .640 & 1.38 & .42 \\
\hline
\end{tabular}
\end{table}

Table~\ref{tab:answers} shows why the lenient gold Score cannot compare
retrievers: it is higher with dense retrieval for both models (e.g.\ 0.809
versus 0.758 for \texttt{qwen2.5:7b}), although dense retrieval supplies fewer
deciding papers for claim questions. The difference comes from accepted
abstentions when the evidence is missing.

With the strict E2E score the overall answer quality is equal. The graph minus
dense difference is $-0.01$ $[-0.10, 0.08]$ for \texttt{qwen2.5:7b}, pooled over
both runs, and $0.00$ $[-0.11, 0.11]$ for \texttt{gpt-oss-20b}. The run-to-run
variance is of the same size as the difference between the systems. Between the
two \texttt{qwen2.5:7b} runs the retrieved papers were identical for 49 of 50
questions with graph retrieval and for all 50 with dense retrieval. Still, the
graph E2E moved from 0.748 to 0.696 and the sign of the difference flipped, and
a single answer moved by 0.08--0.09 on average. Differences of a few points from
a single run should therefore not be read as a difference between systems.

The per-type pattern from retrieval carries over to the answers. On method
comparisons graph retrieval leads in all three measurements (0.60--0.68 versus
0.37--0.51). On claim and comparison questions, answers based on graph
retrieval cite 1.5 to 6 times more grade-2 papers. Verdict accuracy on claim
questions is nevertheless not higher, because the models often decide claims
from prior knowledge rather than from the retrieved evidence. On paper lookup,
dense retrieval leads (0.80 versus 0.61 for \texttt{gpt-oss-20b}).

Answers based on graph retrieval are more faithful to the context in all three
measurements, by 0.02 to 0.05. None of these differences is significant on its
own ($+0.03$ $[-0.03, 0.09]$ for \texttt{qwen2.5:7b}, $+0.05$ $[-0.04, 0.13]$ for
\texttt{gpt-oss-20b}), but the direction is consistent.

\subsection{Nine LLMs on Graph Retrieval}

\begin{table}[t]
\caption{Nine LLMs with graph retrieval (50 questions). Score: gold Score; Test: test split; Uns.: share of unsupported claims; Judge: LLM judge, overall 0--5.}
\label{tab:models}
\centering
\setlength{\tabcolsep}{4pt}
\begin{tabular}{lccccc}
\hline
Model & Score & Test & Faithful & Uns. & Judge \\
\hline
gpt-oss-20b      & \textbf{0.830} & \textbf{0.833} & 0.694 & 0.278 & 4.57 \\
gpt-oss-120b     & 0.798 & 0.745 & 0.722 & 0.246 & 4.53 \\
qwen3.8-27b      & 0.790 & 0.782 & \textbf{0.814} & 0.183 & 4.31 \\
qwen2.5:7b$^\ast$ & 0.758 & 0.779 & 0.771 & 0.218 & 3.46 \\
mistral$^\ast$   & 0.680 & 0.706 & 0.704 & 0.292 & 3.41 \\
llama3$^\ast$    & 0.669 & 0.705 & 0.684 & 0.315 & 3.27 \\
phi3:mini$^\ast$ & 0.663 & 0.656 & 0.703 & 0.284 & 3.49 \\
llama3.2$^\ast$  & 0.648 & 0.676 & 0.668 & 0.332 & 3.39 \\
allam-2-7b       & 0.167 & 0.180 & 0.195 & 0.846 & 1.17 \\
\hline
\multicolumn{6}{l}{\footnotesize $^\ast$ local (Ollama); others via Groq.}
\end{tabular}
\end{table}

Table~\ref{tab:models} reports nine LLMs on the same graph retrieval. With good
retrieval, the best local model, \texttt{qwen2.5:7b}, is within 0.07 of the best
cloud model. On the test split it is above \texttt{gpt-oss-120b} (0.779 versus
0.745), and the larger \texttt{gpt-oss-120b} does not beat \texttt{gpt-oss-20b}.
Given the run-to-run variance measured above (up to 0.05), the order of adjacent
models in Table~\ref{tab:models} is not reliable; the separation between groups
is. The models differ in style: \texttt{qwen3.8-27b} cites nearly all relevant
context papers (CiteUtil 0.98) and is the most faithful, while
\texttt{gpt-oss-20b} selects papers more precisely (CitePrec 0.83) but adds
more content that is not in the context. \texttt{allam-2-7b} is not usable for
the task: it abstains on 78\,\% of answerable questions, and its small context
window causes request failures on some questions.

Contradictions of the context are rare (1--3\,\% of claims, an underestimate
given the judge's behaviour). The main faithfulness problem is addition:
18--33\,\% of claims are not supported by the context. These are mostly
background knowledge, method explanations and results attributed to the wrong
paper.

\subsection{LLM Judge versus Gold Metrics}

The LLM judge agrees with the gold Score only moderately at the answer level
(Spearman $\rho = 0.45$, $n = 450$). At the model level the agreement is high
($\rho = 0.87$), but only because the judge separates cloud from local models;
within the five local models it is $\rho = 0.25$. The judge also overstates the
gap between them: \texttt{gpt-oss-20b} and \texttt{qwen2.5:7b} differ by 0.07 in
gold Score but by 1.1 points (22\,\% of the scale) in the judge's score, and
\texttt{phi3:mini}, second to last among local models by gold Score, is ranked
first among them by the judge. Because the judge sees only source titles, it
cannot reward an answer for citing the right papers or penalise it for
misusing them; it appears to reward style and length instead.

Faithfulness, judged by the same model with the context, measures a
different property from the gold Score: the two correlate at $\rho = 0.67$
across models and only $\rho = 0.15$ across answers. Choosing the right papers
and staying within their content are thus separate abilities, and an
evaluation that reports only one of them misses the other.
